\documentclass[11pt,a4paper]{article}

\usepackage[margin=2.5cm]{geometry}
\usepackage{xcolor}
\usepackage{enumitem}
\usepackage{hyperref}
\usepackage{booktabs}
\usepackage{parskip}

% Formatting for reviewer comments
\definecolor{reviewerblue}{RGB}{0,51,102}
\newenvironment{reviewercomment}{%
  \begin{quote}\itshape\color{reviewerblue}%
}{%
  \end{quote}%
}
\newcommand{\response}{\textbf{Response:}~}

\title{Response to Reviewer Comments\\[0.5em]
\large Manuscript JTB-D-25-00973\\
``Likelihood-Free Parameter Inference for Spatiotemporal\\
Stochastic Biological Models using Neural Posterior Estimation''}
\author{Tom Kimpson, Jennifer Flegg, Matthew J.\ Simpson}
\date{\today}

\begin{document}
\maketitle

\noindent
We thank both reviewers for their detailed and constructive feedback. Their comments have led to substantial improvements in the manuscript. Below we provide a point-by-point response, with all section, table, and figure references referring to the revised manuscript.

\subsection*{Summary of major revisions}

In response to the reviewers' concerns, we have made the following major changes:

\begin{enumerate}[leftmargin=*]
  \item \textbf{IMRaD restructure and biological reframing.} The manuscript has been reorganised into a standard Introduction--Methods--Results--Discussion structure. The narrative now leads with the biological problem (parameter inference for cell migration models) and presents NPE as the tool to address it.
  \item \textbf{Three new models of increasing complexity.} We now present four random walk models: the original model (unbiased diffusion), Model~A (directional bias), Model~B (proliferation), and Model~C (bias + proliferation). This demonstrates NPE across a spectrum from convenience (where classical methods suffice) to necessity (where no tractable likelihood exists).
  \item \textbf{SBI diagnostics.} A complete diagnostic workflow (simulation-based calibration, Tests of Accuracy with Random Points (TARP) coverage tests, classifier two-sample tests) is now reported for all four models (Section~3.5, Table~4, Figure~8).
  \item \textbf{Reproducibility study.} A five-seed reproducibility study demonstrates stable posterior estimates across independent training runs (Appendix~B, Table~7).
  \item \textbf{Timing breakdown.} Computational runtimes are now separated into simulation generation, network training, and inference components for all models (Section~3.6, Table~6).
  \item \textbf{Hyperparameter table.} All key hyperparameters are now reported in the paper (Table~2).
\end{enumerate}


%% ============================================================
\section*{General comments}
%% ============================================================

\noindent
The review form posed nine standard questions to both reviewers. We address each below, quoting both reviewers and providing a unified response.


%% ----- Q1: Objectives -----
\subsection*{Q1 --- Objectives and scope}

\begin{reviewercomment}
\textbf{Are the objectives and the rationale of the study clearly stated?}

\textbf{Reviewer \#1:} Yes, but I have some issues with that. The fundamental issue concerns modelling and simulation in general; its application to biology is just one of several potential applications. As you will see from the paper, the focus is on the method itself rather than its application to biological models. I expand on this topic further in the additional material.

\textbf{Reviewer \#2:} The lack of a clear goal for the study is one of our main concerns. The study provides neither new biological results nor a new method\ldots\ The scope is rather narrow, as the study uses a very specific random walk model that does not strictly require simulation-based inference\ldots\ The study also fails to demonstrate the strengths and challenges of simulation-based inference for a variety of random walk models.
\end{reviewercomment}

\response
We agree that the original manuscript was too method-focused and narrow in scope. The revised manuscript has been restructured into IMRaD format and reframed as a biology paper. The Introduction now opens with the biological problem---parameter inference for lattice-based cell migration models from scratch/barrier assay data---and positions NPE as the inferential tool. The Methods section (Section~2) describes both the biological models and the inference machinery. The Results section (Section~3) presents systematic application across four models of increasing biological complexity, and the Discussion (Section~4) emphasises biological insights: what inferred parameters tell us about cell motility, directional bias, and proliferation. The scope has been substantially broadened with four models (Sections~3.1--3.4). The original model serves as a validation case (where classical methods are available for comparison), while Model~C---combining bias and proliferation with no tractable likelihood---demonstrates where simulation-based inference becomes genuinely necessary.


%% ----- Q2: Replicability -----
\subsection*{Q2 --- Replicability}

\begin{reviewercomment}
\textbf{Are the study methods and analyses replicable?}

\textbf{Reviewer \#1:} Yes.

\textbf{Reviewer \#2:} The code is particularly well written and can be used to reproduce the study\ldots\ There are however some files missing (e.g., \texttt{results\_extracted.pkl} for \texttt{demo.py}) to run notebooks mentioned in the code repository out-of-the-box. It is unclear from the text whether training and parameter inference were performed only once, or if repeating the experiment with different seeds would lead to similar results.
\end{reviewercomment}

\response
We thank both reviewers. The code repository has been updated to include the missing result files needed to run the demonstration notebooks. We have further strengthened replicability by adding a five-seed reproducibility study (Appendix~B, Table~7). Across all four models and all parameters, posterior means are stable (maximum standard deviation of 0.032 across seeds) and 95\% credible intervals achieve 100\% coverage of the true values.


%% ----- Q3: Statistical analyses -----
\subsection*{Q3 --- Statistical analyses}

\begin{reviewercomment}
\textbf{Are the statistical analyses and their interpretation appropriate?}

\textbf{Reviewer \#1:} Yes.

\textbf{Reviewer \#2:} The analysis does not report diagnostic of the simulation-based inference (although the authors provide posterior predictive checks, which are important, they are not a complete diagnostic workflow)---which is not a problem here because parameters were known from another study, but which would, in general, be a problem. Reference~15 of the article provides guidelines for such diagnostics.
\end{reviewercomment}

\response
For completeness we have added a complete SBI diagnostic workflow following Cranmer et al.\ (2020), now reported in Section~3.5. For each of the four models, we perform:
\begin{itemize}[nosep]
  \item \textbf{Simulation-based calibration (SBC)}: 1000 simulations with rank uniformity assessed via Kolmogorov--Smirnov tests. SBC rank empirical cumulative distribution function (ECDF) plots are shown in Figure~8.
  \item \textbf{Classifier two-sample tests (C2ST)}: Scores near 0.5 across all models confirm well-learned posteriors.
  \item \textbf{TARP coverage tests}: Area-under-the-curve (AUC) values quantify calibration quality.
\end{itemize}
Results are summarised in Table~4.


%% ----- Q4: Figures and tables -----
\subsection*{Q4 --- Figures and tables}

\begin{reviewercomment}
\textbf{Are the figures and tables appropriate?}

\textbf{Reviewer \#1:} Yes, because I have the impression that the scope of the journal is not being met.

\textbf{Reviewer \#2:} Section~4.3.2 states: ``Figure~6 shows the joint posterior distribution inferred from the 2D spatial data, overlaid with the result from the 1D column count analysis for direct comparison.'' However, this does not seem to be the case for Figure~6. It would be helpful to include a summary table of the hyperparameters in the paper itself, rather than only in the code repository.
\end{reviewercomment}

\response
We have substantially revised the figures and tables. All posterior figures now use separate subfigures for the 1D and 2D results (Figures~4--7), resolving the overlay issue in the original Figure~6. The misleading text has been removed. A hyperparameter summary table has been added (Table~2). The timing table has been expanded to separate simulation, training, and inference components for all four models (Table~6). SBC diagnostic plots have been added (Figure~8). Posterior predictive checks are presented in Appendix~A.


%% ----- Q5: Interpretations and conclusions -----
\subsection*{Q5 --- Interpretations and conclusions}

\begin{reviewercomment}
\textbf{Are the interpretations and conclusions justified by the results?}

\textbf{Reviewer \#1:} The problem is that only one example was chosen to demonstrate the practical application of the method. For such a topic, I would expect this method to be applied to a series of examples, including a comparison and benchmarking.

\textbf{Reviewer \#2:} N/A.
\end{reviewercomment}

\response
The revised manuscript now presents four random walk models of increasing biological complexity:
\begin{itemize}[nosep]
  \item \textbf{Original model} (Section~3.1): unbiased diffusion with two parameters ($U$, $P$), benchmarked against ABC, surrogate-likelihood MLE/Laplace, and surrogate-likelihood MCMC (Table~3, Figure~3).
  \item \textbf{Model~A} (Section~3.2): directional bias with three parameters ($U$, $P$, $\rho$).
  \item \textbf{Model~B} (Section~3.3): proliferation with three parameters ($U$, $P$, $R$).
  \item \textbf{Model~C} (Section~3.4): combined bias and proliferation with four parameters ($U$, $P$, $\rho$, $R$), where no tractable likelihood exists.
\end{itemize}
Comprehensive NPE results across all models are collected in Table~5. The progression from the original model (where classical methods work well) to Model~C (where they are unavailable) demonstrates when NPE transitions from a convenience to a necessity.


%% ----- Q6: Strengths -----
\subsection*{Q6 --- Strengths}

\begin{reviewercomment}
\textbf{Have the authors clearly emphasized the strengths of their study/theory/methods/argument?}

\textbf{Reviewer \#1:} Yes.

\textbf{Reviewer \#2:} Because of the lack of clearly established goals, it is currently unclear what would be the current strengths of the manuscript.
\end{reviewercomment}

\response
We have clarified the goals of the study through the biological reframing described above. The strengths---amortized inference, scalability to intractable likelihoods, and automated feature extraction via CNNs---are now tied explicitly to the biological application.


%% ----- Q7: Limitations -----
\subsection*{Q7 --- Limitations}

\begin{reviewercomment}
\textbf{Have the authors clearly stated the limitations of their study/theory/methods/argument?}

\textbf{Reviewer \#1:} Not really! For such a topic, I would expect this method to be applied to a series of examples, including a comparison and benchmarking.

\textbf{Reviewer \#2:} It is not clear why authors chose NPE specifically? What about other algorithms? (e.g., NRE)---this choice should probably be at least discussed or justified. [The 1D vs 2D comparison] does not seem to illustrate an improvement in the precision of parameter inference. It would be much more convincing to present a case in which switching from 1D to 2D summary statistics actually impacts the quality of the parameter estimate.
\end{reviewercomment}

\response
As described in Q5, we now present four models with benchmarking against classical methods.

Regarding NPE vs alternatives: we have added a discussion in Section~4.1. The key considerations are: (i)~NPE learns the posterior directly, providing immediate samples for any new observation without additional MCMC sampling (unlike NRE, which learns likelihood-to-evidence ratios and requires MCMC to obtain the posterior); (ii)~NPE's amortized nature is particularly advantageous when posteriors are needed for many observations, as in experimental screening or time-course analysis; and (iii)~NPE has a strong track record in scientific inverse problems. We cite relevant benchmark comparisons (Lueckmann et al.\ 2021) and note that the choice between NPE and NRE is ultimately problem-dependent.

Regarding the 1D vs 2D comparison: we applied the 2D CNN approach to Model~A (directional bias), where spatial anisotropy carries information that 1D column counts collapse by summing across rows. As shown in Section~3.2 and Figure~5, the 2D posteriors demonstrate the CNN's ability to recover this spatial structure. We have reframed the comparison in Section~4.2: the 1D column counts are hand-crafted summary statistics designed with domain knowledge, whereas the 2D CNN performs automated feature extraction directly from raw spatial data. The primary advantage of the 2D approach is thus not necessarily tighter posteriors in all cases, but rather the automation of feature engineering---a critical capability when designing informative summary statistics is non-trivial or when the relevant spatial features are unknown \emph{a priori}.


%% ----- Q8: Structure and writing -----
\subsection*{Q8 --- Structure and writing}

\begin{reviewercomment}
\textbf{Is the manuscript well-structured and clearly written?}

\textbf{Reviewer \#1:} [In the additional material:] Adopt IMRaD format; clarify Discussion vs Conclusion.

\textbf{Reviewer \#2:} The writing could benefit, at times, from more precision and fewer colloquialisms. For example\ldots\ ``deceptively simple'' or ``familiar inverse transform sampling'' could be avoided and enhanced.
\end{reviewercomment}

\response
The manuscript has been restructured into IMRaD format, and the Discussion and Conclusion have been merged into a single Discussion section (Section~4) with subsections on: advantages and limitations of NPE (Section~4.1), the role of spatial information (Section~4.2), and computational scalability and future directions (Section~4.3). We have also removed the identified colloquialisms (``deceptively simple'', ``familiar'') and performed a general pass for imprecise language throughout the manuscript.


%% ----- Q9: Language editing -----
\subsection*{Q9 --- Language editing}

\begin{reviewercomment}
\textbf{Does the language need editing?}

\textbf{Reviewer \#1:} No.

\textbf{Reviewer \#2:} No.
\end{reviewercomment}

\response
We appreciate the reviewers' assessment. We have nonetheless performed a careful language editing pass to improve precision and readability throughout the manuscript.


%% ============================================================
\section*{Specific comments}
%% ============================================================


%% ----- Reviewer #1 specific -----
\subsection*{Reviewer \#1 --- Annotated PDF}

\begin{reviewercomment}
Various annotations in the attached PDF regarding: abbreviations in abstract, formula integration (colons before equations), Figure~2 captions, Figure~5 placement, and minor typos.
\end{reviewercomment}

\response
All items from the annotated PDF have been addressed:
\begin{itemize}[nosep]
  \item \textbf{Abbreviations in abstract}: All abbreviations (NPE, CNN) have been spelled out in full in the abstract.
  \item \textbf{Formula integration}: Colons preceding equations have been removed at six locations where the sentence continues after the formula.
  \item \textbf{Figure~2 captions}: Subfigure descriptions have been moved into the main caption body as (a)\ldots, (b)\ldots, etc.
  \item \textbf{Figure placement}: Figure placement has been reviewed and corrected following the IMRaD restructure.
  \item \textbf{Typos}: All highlighted typos have been corrected, including hyphen-to-en-dash fixes, notation inconsistencies, and grammar corrections.
\end{itemize}


%% ----- Reviewer #2 specific -----
\subsection*{Reviewer \#2 --- Additional comments}

\begin{reviewercomment}
The manuscript would benefit from a clearer direction with either more novelty in the research results or more general guidance for using simulation-based inference for biological models. A suggestion for such a direction could be, for instance, to explore the strengths and weaknesses of using simulation-based inference methods for biological random-walk models in general.
\end{reviewercomment}

\response
We have followed this suggestion closely. The revised manuscript now explores NPE across four random walk models of increasing complexity, systematically demonstrating its strengths (amortized inference, scalability to intractable likelihoods, automated feature extraction) and challenges (calibration of proliferation parameters, the simulation bottleneck for complex models). The Discussion (Section~4) explicitly addresses when NPE is a convenience versus a necessity for biological modelling.


%% ----- R2 item 2: Amortized definition -----
\subsubsection*{Item 2 --- Definition of ``amortized''}

\begin{reviewercomment}
P7: the definition of ``amortized'' could come earlier for clarity, the word being mentioned twice earlier (p3, in Introduction 1 and Table 1).
\end{reviewercomment}

\response
The definition has been moved to its first use in the Introduction: ``\ldots amortized inference (where the upfront cost of training is offset by near-instant posterior estimation for any new observation)\ldots''


%% ----- R2 item 3: SBI package reference -----
\subsubsection*{Item 3 --- SBI package reference}

\begin{reviewercomment}
As stated in the SBI package documentation, reference~63 (SBI package) should be updated to the JOSS publication.
\end{reviewercomment}

\response
The reference has been updated to Boelts et al.\ (2025), \emph{Journal of Open Source Software}, 10(108), 7754.


%% ----- R2 item 4: Section 3.2.1 unclear -----
\subsubsection*{Item 4 --- Section 3.2.1 unclear}

\begin{reviewercomment}
Section~3.2.1 does not seem very clear.
\end{reviewercomment}

\response
This section has been substantially rewritten as part of the IMRaD restructure. The former Section~3.2.1 content is now integrated into the Methods section (Section~2), where the inference pipeline is described in a coherent, self-contained manner.


%% ----- R2 item 5: SMC-ABC mention -----
\subsubsection*{Item 5 --- SMC-ABC description}

\begin{reviewercomment}
The last paragraph of Section~2.1: ``More efficient variants of ABC, such as Sequential Monte Carlo ABC, have been proposed\ldots'' It is unclear why the authors mention alternative methods without discussing them.
\end{reviewercomment}

\response
The description of SMC-ABC has been expanded to explain the key idea: SMC-ABC iteratively refines proposal distributions and progressively lowers the acceptance tolerance, improving sampling efficiency relative to rejection-based ABC. This is now in the Classical Inference Approaches subsection of Section~2 (Section~2.2).


%% ----- R2 item 6: Section 2.3 label -----
\subsubsection*{Item 6 --- Section 2.3 label ``Example Results''}

\begin{reviewercomment}
Section~2.3 label ``Example Results'' is unclear.
\end{reviewercomment}

\response
The IMRaD restructure resolves all section naming issues. The former ``Example Results'' is now part of the Results section (Section~3).


%% ----- R2 item 7: ``Fundamentally different'' -----
\subsubsection*{Item 7 --- ``Fundamentally different'' phrasing}

\begin{reviewercomment}
``NPE offers a fundamentally different approach to inference for simulator-based models.'' This statement is incorrect or unclear. NPE is a simulation-based inference method that produces approximation of a posterior distribution over parameters.
\end{reviewercomment}

\response
We agree that this phrasing was imprecise. The text now reads ``distinct approach'', accurately positioning NPE as one approach within the broader SBI framework rather than a fundamentally different paradigm.


%% ----- R2 item 8: SNPE / Neural Spline Flows -----
\subsubsection*{Item 8 --- SNPE confusion and Neural Spline Flows}

\begin{reviewercomment}
Section~3.3 describes SNPE but the closing sentence says ``we will exclusively use the sequential version of NPE''. Neural Spline Flows are not detailed even though they are used.
\end{reviewercomment}

\response
The SNPE/NPE confusion has been resolved: we now clarify that single-round NPE is used throughout (Section~2.3, Implementation Details). Neural Spline Flows are now described in the Implementation Details subsection of Section~2.3, including their use of monotonic rational-quadratic spline transformations with analytic Jacobians.


%% ----- R2 item 9: ``Minimally informative'' prior -----
\subsubsection*{Item 9 --- ``Minimally informative'' prior}

\begin{reviewercomment}
``In the absence of strong prior information about parameter values, these uniform priors represent a minimally informative choice that allows the data to dominate the posterior inference.'' We believe this statement is imprecise, as uniform priors are not necessarily the minimally informative choice in general.
\end{reviewercomment}

\response
We agree. The imprecise phrasing has been removed; the revised manuscript instead describes the priors as ``broad uniform priors determined by the physical constraints of the model parameters'', avoiding any claim about minimal informativeness.


%% ----- R2 item 10: Table 1 -----
\subsubsection*{Item 10 --- Table~1: Cost per Inference vs Amortized}

\begin{reviewercomment}
Cost per Inference vs.\ Amortized---are they the same thing?
\end{reviewercomment}

\response
These are distinct concepts. ``Amortized'' refers to the property of the method (a one-time training cost enables inference for any new observation), while ``cost per inference'' is the resulting computational cost for each posterior estimate. We have clarified this by renaming the column to ``Per-observation Cost'' and adding a footnote to the ``Amortized'' column explaining the distinction (Table~1).


%% ----- R2 item 11: Comparative results table -----
\subsubsection*{Item 11 --- Comparative results table}

\begin{reviewercomment}
It would be great to include the results from the alternative models (ABC, MCMC, etc.) and the true values.
\end{reviewercomment}

\response
Table~3 now includes the true parameter values and results from three classical methods: ABC (with SMC), surrogate-likelihood MLE with Laplace approximation, and surrogate-likelihood MCMC. The NPE results are shown alongside for direct comparison. Comprehensive results across all four models are collected in Table~5.


%% ----- R2 item 12: Timing table -----
\subsubsection*{Item 12 --- Timing breakdown}

\begin{reviewercomment}
It would be informative to separate the training of the network and the generation of the simulation for training, since simulations are usually the bottleneck for complex simulators rather than for training the network.
\end{reviewercomment}

\response
Table~6 now separates computational costs into three components---simulation generation, neural network training, and per-observation inference---for all four models in both 1D and 2D configurations. This reveals that simulation generation dominates the computational budget: it accounts for 74\% of the total cost for the original model and over 96\% for the exclusion-process models (A--C), where the per-simulation cost increases from $\sim$20 sims/s to $\sim$2 sims/s due to the exclusion constraint. Neural network training on an A100 GPU requires only 2--3 minutes in all cases.


\end{document}
